%%%%%%%%%%%%
%
% $Autor: Wings $
% $Datum: 2019-03-05 08:03:15Z $
% $Version: 4250 $
% !TeX spellcheck = en_GB/de_DE
% !TeX encoding = utf8
% !TeX TXS-program:bibliography = txs:///biber
%
%%%%%%%%%%%%

\chapter{Object Recognition with Jetson Nano} 

\section{Solution Architecture for Object Detection with JetsonNano}
For this study in the project to detect and classify the objects, the \ac{kdd} process is used.The \ac{kdd} is the process of discovering useful knowledge from a collection of data. This technique includes steps such as data selection, preparation, using right algorithms etc.,
\begin{figure}[H]
	\centering
	
	\GRAPHICS{0.8}{0.6}{CreatingDatabase/SolutionArchitecture.png}
	
	\caption{Solution Architecture for Object Detection using JetsonNano}
\end{figure} 

\section{Description of the Objects}

The task is to detect 2 types of objects: a cuboid and a cylinder.The objects are made of metal and dark in color.
The dimensions of the cylindrical object are: length: 11 cm and diameter: 4cm
\begin{figure}[H]
	\centering
	
	\GRAPHICS{1.0}{1.0}{CreatingDatabase/CylinderMeasure.png}
	
	\caption{The dimensions of the cylindrical object}
\end{figure} 
The dimensions of the cuboid object are: length: 11 cm, width: 1 cm and breadth: 5 cm.
\begin{figure}[H]
	\centering
	
	\GRAPHICS{1.0}{1.0}{CreatingDatabase/CuboidMeasure.png}
	
	\caption{The dimensions of the cuboid object}
\end{figure} 

\section{Database}

The images are captured using the Raspberry PI camera v2.1 attached to the Jetson Nano.

The code uses the function \PYTHON{gstreamer\_pipeline()}. GStreamer is a pipeline-based multimedia framework that links together a wide variety of media processing systems to complete complex workflows. For instance, GStreamer can be used to build a system that reads files in one format, processes them, and exports them in another. The formats and processes can be changed in a plug and play fashion.GStreamer supports a wide variety of media-handling components, including simple audio playback, audio and video playback, recording, streaming and editing. The pipeline design serves as a base to create many types of multimedia applications such as video editors, transcoders, streaming media broadcasters and media players.

Frames were captured from the livestream by the JetsonNano using a python program. The frame rate is defined as 5fps and the dimensions of each frame is 1280 $\times$ 720. The program is written to sequentially store the captured frames. The module cv2 is used for this. Frames are read using the function  \PYTHON{cv2.videoCapture()} where we can specify flip values to flip the frame upside down or capture the stream as it is.The flip\_method is set as 2.There are 2 values for the frame\_method which are 1 and 2. If the flipmethod is 0 then the frame is inverted.  The frames are then written in a folder using the function\PYTHON{cv2.imwrite(path, image)} where path is the complete path of the output file to which we would like to write the image numpy array.The function returns a boolean value. True if the image is successfully written and False if the image is not written successfully to the local path specified.
The below is the result of running the program to capture frames through video stream in the terminal of Jetson Nano.

\begin{figure}[H]
	\centering
	
	\GRAPHICS{0.8}{0.8}{CreatingDatabase/SimpleCameraRunning.png}
	
	\caption{Result of running the program to capture frames through video stream in Jetson Nano}
\end{figure}  
The below is the python code used: 
\begin{code}[H]
	\lstinputlisting[language=Python, firstnumber=1]{../Code/JetsonNano/SimpleCamera.py}
	\caption{Program to capture the frames from the camera stream}\label{TensorFlow:Complete}
\end{code}

\subsection{Sample images from the dataset for cuboid object}

For the cuboid, around 1352 images were created in different angles in approx. 7 minutes with a frame rate of 5 \ac{fps}.Below are the sample set of images from the dataset of the object cuboid in various positions. The images are captured in whitebackground to improve the accuracy of the machine learning algorithm. 
Below are the sample set of images from the dataset of the object cuboid in various positions.The frames were stored sequentially in the \FILE{.jpg} format in the folder specified. 
\begin{figure}[H]
	\centering
	\GRAPHICS{0.8}{0.8}{CreatingDatabase/CuboidSamples.png}
	\caption{Sample images for cuboid captured in different positions}
\end{figure} 
\subsection{Sample images from the dataset for the cylindrical object}
For the cylindrical object around 1126 images were captured in different angles with a white background in around approximately 7 minutes with a frame rate of 5 \ac{fps}.
Below are the sample set of images from the dataset of the object cylinder in various positions. 
\begin{figure}[H]
	\centering
	\GRAPHICS{0.8}{0.8}{CreatingDatabase/CylinderSamples.png}
	\caption{sample images for cylinder captured in different positions}
\end{figure} 

\subsection{Validation of the dataset}

Below are the observations on the dataset prepared:
\begin{itemize}
	\item Poor lighting conditions.The lighting conditions should be good to achieve a good performance by the Machine learning algorithm
	\item Poor sharpness of the image.
	\item Bad background in some images.
\end{itemize}

\subsection{Possible improvements}

Below are some of the improvements that can be done: 
\begin{itemize}
	\item Calibrate the lens of the raspberry PI camera. 
	\item Re-capture the images with good background. 
	\item Introducing a wait time between frames captured to avoid the repetition of the images. 
	\item Replace the camera with a high quality cam.  
\end{itemize}

\subsection{Possible Re-work if needed }

\begin{itemize}
	\item Preparing the database again with a better background and lighting. 
	\item Study about the methods to improve the lighting captured by the camera of Jetson Nano 
\end{itemize}
Images were captured again with a better background and using LED lighting. Frame rate is set at 5 \ac{fps}. No significant improvement observed. Improper images with blurr and bad background and lighting were removed. With the remaining images,a dataset is prepared for the 2 classes: cylindrical and cuboid objects from the captured images. Details about the dataset preparation are mentioned in the below section.
